\documentclass[11pt]{article}
\usepackage[margin=25mm]{geometry}
\usepackage{fontspec,unicode-math}
\setmainfont{Latin Modern Roman}
\setmathfont{Latin Modern Math}
\usepackage{amsmath,amssymb,mathtools}
\usepackage{xurl,hyperref}
\hypersetup{colorlinks=true,urlcolor=blue,linkcolor=blue}
\setcounter{secnumdepth}{0}
\setlength{\parindent}{0pt}
\setlength{\parskip}{5pt}
\setlength{\emergencystretch}{3em}
\newcommand{\sref}[2]{\hyperlink{source:#1:#2}{[S#2]}}
\newcommand{\cataloguescope}[1]{\paragraph{Source question.}#1}
\newcommand{\floor}[1]{\left\lfloor#1\right\rfloor}
\begin{document}
\section{Duality and Convergence Rates in Lasserre Hierarchies}\label{statement}

\subsection{Definitions and mathematical statement}

\paragraph{Shared notation.}Standard number systems are denoted by $\mathbb N,\mathbb Z,\mathbb Q,\mathbb R,\mathbb C$; index ranges, computational models and specialized notation follow the cited problem source and the definitions below.


For a polynomial optimization instance $\inf_{x\in K}f(x)$ with $K=\{x:g_j(x)\geq0\}$, a representing measure $\mu$ has moments $y_\alpha=\int_K x^\alpha\,d\mu(x)$. Moment relaxations retain finitely many moments and require moment and localizing matrices to be positive semidefinite. Their SOS duals seek certificates of the form $f-\gamma=\sigma_0+\sum_j\sigma_jg_j$, with sum-of-squares polynomials and bounded degrees. The generalized moment problem also permits the source linear measure constraints. Strong duality means equality of the primal and dual optimal values under the appropriate attainment conventions. Convergence of moments to an optimal measure and convergence rates of relaxation values are separate targets.
\paragraph{Statement recorded in the supplied source.}
For Lasserre's hierarchy, when does strong duality hold, do the moment relaxations converge to an optimal measure, and what is their convergence rate? \sref{statement}{2}
\paragraph{Scope qualification.}
Compactness, Archimedean certificates and constraint qualifications are hypotheses to be investigated, not silently imposed. Multiple optimal measures can prevent a full moment sequence from converging to one measure without an additional selection rule.
\paragraph{Source status.}
The supplied catalogue labels this record partially solved. Its source-status metadata was checked against the supplied snapshot on 2026-10-11. This import does not independently establish current open status or mathematical correctness. The catalogue literature assessment is dated 2026-08-22.
\subsection{Short English statement}
Identify when moment and SOS bounds have no duality gap, when optimizing moments converge to an optimal measure, and how fast the bounds converge.
\subsection{Sources}
\begin{itemize}
\item[\hypertarget{source:statement:1}{S1}] Ulam AI and the UnsolvedMath Contributors, UnsolvedMath: A Curated Collection of Open Mathematics Problems, catalogue entry. Catalogue curation and metadata: CC BY 4.0. The supplied source question is retained; definitions and the English summary are editorial additions.
\url{https://huggingface.co/datasets/ulamai/UnsolvedMath}.
\item[\hypertarget{source:statement:2}{S2}] Mixed-integer Nonlinear Optimization: A Hatchery for Modern Mathematics (2024), Oberwolfach Reports, DOI 10.4171/owr/2023/35, pp. 1953--2016
\url{https://doi.org/10.4171/owr/2023/35}.
\item[\hypertarget{source:statement:3}{S3}] Tacchi
\url{https://arxiv.org/abs/2011.08139}.
\item[\hypertarget{source:statement:4}{S4}] Josz-Henrion
\url{https://arxiv.org/abs/1405.7334}.
\item[\hypertarget{source:statement:5}{S5}] Kirschner-de Klerk
\url{https://doi.org/10.1007/s11590-022-01851-3}.
\textit{Underlying Oberwolfach Reports material retains its source-specific licence, including CC BY-SA 4.0 where applicable; see THIRD\_PARTY\_NOTICES.md.}
\end{itemize}
\paragraph{Preparation and review.}
Original question attributed to the cited source and UnsolvedMath contributors. Codex (GPT-6; exact serving revision not exposed) assisted with definitions, English summaries, source inspection and formatting. No new human mathematical authorship or independent proof verification is asserted. Catalogue identity checked against the complete supplied registry and existing repository statements; known flagged repetitions and obvious identical conjectures excluded. Full mathematical-equivalence/literature review has not been performed. Statement preparation and final review followed contribution prompts C and D; automated checks verify structure and identity, not open status.
\label{end:statement}
\end{document}
